% Source: 05 - Tue Oct 6 2026.pdf, page 3. Content remains in source order.
\sourcepage{05}{3}
\textbf{OBSERVATION:} $\psi(\vec x,\vec y)$ is always satisfiable (HOMEWORK: how many satisfying assignments?)

We have that
\[|\langle\psi(\vec x,\vec y)\rangle|=\Theta(|\langle C(\vec x)\rangle|)\]
since each subexpression uses a number of symbols linear in the size of the corresponding gate!

The reduction function $f$ uses $\psi$ as follows:
\[
f(\langle C(\vec x)\rangle)=\phi_C(\vec x,\vec y)=\psi(\vec x,\vec y)\land\underbrace{y_0}_{\text{variable associated to gate of output wire}}.
\]
\[
|\langle\phi_C(x,y)\rangle|=\Theta(|\langle\psi(\vec x,\vec y)\rangle|)=\Theta(|\langle C(x)\rangle|).
\]
Thus $f$ is ptc.
